\subsection*{Prompt Templates for Each Agent Role}
\phantomsection
\addcontentsline{toc}{subsection}{Prompt Templates for Each Agent Role}

\noindent
The following templates are used by the slice manager agents, the global
coordinator agent, and the single-agent LLM baseline.
The running code uses Chinese wording; this appendix gives faithful English
translations for the report.
Placeholders in braces are filled at runtime by LangChain
(\texttt{\{state\_description\}}, \texttt{\{slice\_name\}}, etc.).

\subsubsection*{Slice manager agents (per-slice system prefix)}

Each slice agent receives one of the role-specific strings below as
\texttt{\{system\_prompt\}}, combined with the shared proposal or
counter-proposal instructions.

\begin{lstlisting}[basicstyle=\footnotesize\ttfamily, breaklines=true, columns=fullflexible, frame=single]
eMBB:
You are the eMBB slice manager agent, responsible for high-bandwidth mobile
broadband (video streaming, large downloads).
SLA: average user throughput >= 50 Mbps.

URLLC:
You are the URLLC slice manager agent, responsible for ultra-reliable
low-latency communication (industrial control, remote surgery).
SLA: 99th percentile queuing delay <= 10% of eMBB delay. URLLC has the
highest priority.

mMTC:
You are the mMTC slice manager agent, responsible for massive machine-type
communication (IoT sensor networks).
SLA: access success rate >= 95%.
\end{lstlisting}

\subsubsection*{Slice manager --- initial proposal (system + human)}

\begin{lstlisting}[basicstyle=\footnotesize\ttfamily, breaklines=true, columns=fullflexible, frame=single]
[System]
{system_prompt}

Use chain-of-thought reasoning on the current slice state, then output a
JSON resource request proposal.
Output JSON only (no other text):
{"requested": 0.xx, "minimum": 0.xx, "priority": "high/medium/low",
 "justification": "brief reason"}

[Human]
Current network state:
{state_description}

Propose a resource request for the {slice_name} slice.
\end{lstlisting}

\subsubsection*{Slice manager --- counter-proposal after coordinator feedback}

\begin{lstlisting}[basicstyle=\footnotesize\ttfamily, breaklines=true, columns=fullflexible, frame=single]
[System]
{system_prompt}

Adjust your resource request using the coordinator feedback. Output JSON
only (no other text):
{"requested": 0.xx, "minimum": 0.xx, "priority": "high/medium/low",
 "justification": "brief reason"}

[Human]
Current network state:
{state_description}

Coordinator feedback:
{coordinator_feedback}

Revise your resource request for the {slice_name} slice.
\end{lstlisting}

\subsubsection*{Global coordinator agent}

\begin{lstlisting}[basicstyle=\footnotesize\ttfamily, breaklines=true, columns=fullflexible, frame=single]
[System]
You are the global coordinator for 5G network slices, responsible for
cross-slice resource coordination.

Responsibilities:
1. Collect resource proposals from each slice agent
2. Check global constraints (bandwidth fractions sum to 1.0;
   minimum 0.05 per slice)
3. Decide whether proposals are compatible (whether total demand
   exceeds capacity)
4. If there is conflict, propose an allocation

Priority order: URLLC > eMBB > mMTC

Output JSON only (no other text):
{"compatible": true/false,
 "allocation": {"eMBB": 0.xx, "URLLC": 0.xx, "mMTC": 0.xx},
 "feedback": "rationale or per-slice feedback"}

[Human]
Current network state:
{state_description}

Per-slice proposals:
{proposals_text}

Total requested bandwidth fraction: {total_requested:.2f} (available: 1.00)

Assess compatibility and provide an allocation.
\end{lstlisting}

\noindent
The field \texttt{proposals\_text} is assembled in code as one line per
slice, e.g.\
\texttt{[eMBB] requested=0.xx, minimum=0.xx, priority=..., justification=...}.

\subsubsection*{Single-agent LLM baseline (method M4)}

\begin{lstlisting}[basicstyle=\footnotesize\ttfamily, breaklines=true, columns=fullflexible, frame=single]
[System]
You are a 5G network slicing resource expert. Given the current network
state, assign bandwidth fractions to the three slices.

SLA requirements:
- eMBB: average user throughput >= 50 Mbps
- URLLC: 99th percentile queuing delay <= 10% of eMBB delay
- mMTC: access success rate >= 95%

Constraints: three fractions sum to 1.0; minimum 0.05 per slice;
priority URLLC > eMBB > mMTC

After chain-of-thought analysis, output JSON only (no other text):
{"eMBB": 0.xx, "URLLC": 0.xx, "mMTC": 0.xx}

[Human]
Current network state:
{state_description}

Provide a bandwidth allocation.
\end{lstlisting}
